\documentclass[a4paper,12pt]{extarticle}
\usepackage[table,xcdraw]{xcolor}
\usepackage{pgfplots}
\usepackage[unicode, pdftex]{hyperref}
\usepackage[T2A]{fontenc}
\usepackage[utf8]{inputenc}
\usepackage[russian, english]{babel}
\usepackage{amsmath,amsfonts,amsthm, mathtools}
\usepackage{amssymb}
\usepackage{graphicx}
\graphicspath{{images/}}
\usepackage{wrapfig}
\usepackage{svg}
\usepackage{floatrow}
\usepackage{setspace}
\usepackage[headings]{fancyhdr}
\usepackage[margin=10pt,font={small,stretch=0,9},labelfont=bf,labelsep=period,
justification=centerlast]{caption}
\usepackage{array,tabularx,tabulary,booktabs}
\RequirePackage{multirow}
\RequirePackage{multicol}
\RequirePackage{longtable}
\usepackage{indentfirst}
\usepackage{hyperref}
\usepackage{icomma}
\newcommand*{\hm}[1]{#1\nobreak\discretionary{}
	{\hbox{$\mathsurround=0pt #1$}}{}}
\usepackage{soulutf8} 
\usepackage{etoolbox} 
\usepackage{geometry}
\geometry{top=20mm}
\geometry{bottom=20mm}
\geometry{left=20mm}
\geometry{right=20mm}
\DeclareMathOperator{\sgn}{\mathop{sgn}}
\usepackage[OT1]{fontenc}
\usepackage{amsmath}
\usepackage{amsfonts}
\usepackage{amssymb}
\usepackage{wasysym}
\usepackage{wrapfig}
\usepackage{physics}
\usepackage[normalem]{ulem}
\graphicspath{{pictures/}}
\DeclareGraphicsExtensions{.png,.jpg,.jpeg}
\usepackage{indentfirst}
\usepackage[T2A]{fontenc}
\usepackage{subfigure}
\usepackage{fancyhdr}
\usepackage{caption}

\usepackage{tabularx}
\usepackage{floatrow}
\usepackage{multicol}
\usepackage{lastpage}

\pgfplotsset{
  MyStyle1/.style={
    axis lines = middle,
    enlargelimits = true,
    grid=both,
    grid style={line width=.1pt, draw=gray!10},
    major grid style={line width=.2pt,draw=gray!50},
    minor tick num=5,
    enlargelimits=0.05,
    ticklabel style={font=\tiny,fill=white},
    xlabel style={at={(ticklabel* cs:1)},anchor=north west},
    ylabel style={at={(ticklabel* cs:1)},anchor=south east},
  }
}
\def\Slope{1.5}
\def\Intercept{-20}


% Настройка отчёркиваний и нумерации
\usepackage{fancyhdr}
\pagestyle{fancy}
\fancyhf{}

\fancyhead[L]{\nouppercase{\leftmark\hfill}}
\fancyfoot[RO, RE]{Страница \thepage \ из \pageref{LastPage}}
\renewcommand{\headrulewidth}{0.5pt}
\renewcommand{\footrulewidth}{0.5pt}

\usepackage{titlesec}
\titlelabel{\thetitle.\quad}

\pgfplotsset{compat=1.18}
\begin{document}

\thispagestyle{empty}
\begin{center}
	\textit{Федеральное государственное автономное образовательное\\ учреждение высшего образования }
	
	\vspace{0.5ex}
	
	\textbf{«Московский физико-технический институт\\ (национальный исследовательский университет)»}
\end{center}

\vspace{10ex}
\begin{figure}[!h]
  \begin{center}
    
\includegraphics[width = 0.8\textwidth]{MIPT.png}
  \end{center}
\end{figure}

\begin{center}
	
	\textbf{Лабораторная работа № 3.4.5}
	
	\vspace{1ex}
	
	по курсу общей физики
	
	на тему:
	
	\textbf{\textit{<<Петля гистерезиса>>}}
	
	\vspace{20ex}
	
	\begin{flushright}
		\noindent
		\textit{Работу выполнил:}\\  
		\textit{Легких Никита \\(группа Б05-301)}
	\end{flushright}
	\vfill
	Долгопрудный \\ \today
	\setcounter{page}{1}
\end{center}
\newpage


\setcounter{page}{2} %это нумерация страницы
\let \savenumberline \numberline
\def \numberline#1{\savenumberline{#1.}}
\bigskip

% Конец шаблона
\newpage

\textbf{Цель работы:} изучение петель гистерезиса различных ферромагнитных
материалов в переменных полях.

\textbf{В работе используются:} автотрансформатор, понижающий трансформатор, интегрирующая цепочка, амперметр, вольтметр, электронный
осциллограф, делитель напряжения, тороидальные образцы с двумя обмотками.


\section*{Теоретические сведения}

Основные характеристики
ферромагнетиков — их коэрцитивное поле $H_c$, магнитная проницаемость
$\mu$, рассеиваемая в виде тепла при перемагничивании мощность — зависят
от частоты перемагничивающего поля. В данной работе кривые гистерезиса ферромагнитных материалов изучаются в поле частоты $\nu_0 =$ 50 Гц
с помощью электронного осциллографа.

\begin{figure}[H]
	\centering
	\includegraphics[width=0.4\linewidth]{3.4.5 Петля гистерезиса/pictures/gist3.jpg}
	\caption{Петля гистерезиса ферромагнетика}
	\label{fig:petlya}
\end{figure}

Магнитная индукция $ B $ и напряжённость поля $ H $ в ферромагнитном материале неоднозначно связаны между собой: индукция зависит
не только от напряжённости, но и от предыстории образца. Связь между $ B $ и $ H $ типичного ферромагнетика иллюстрирует рис. $\ref{fig:petlya}$.

Если к ферромагнитному образцу прикладывать переменное внешнее
магнитное поле, то его состояние на плоскости $ B-H $ будет изменяться
по замкнутой кривой — петле гистерезиса. Размер петли определяется
максимальным значением напряжённости $ H $ в цикле (например, петля $ AA' $,
обозначенная пунктиром на рис. $\ref{fig:petlya}$). Если амплитуда напряжённости достаточно велика, то образец будет периодически достигать насыщения,
что на рисунке соответствует кривой $ CEFC'E'F'C $ (предельная петля
гистерезиса). Пересечение предельной петли с вертикальной осью соответствует остаточной индукции $B_r$, пересечение с горизонтальной осью
— коэрцитивному полю $H_c$. Крайние точки петель, соответствующие амплитудным значениям $ H $ (например, точка $ A $ на рис. $\ref{fig:petlya}$), лежат на начальной кривой намагничивания ($ OAC $).

\textbf{Измерение магнитной индукции.} Магнитную индукцию $ B $ удобно
определять с помощью ЭДС, возникающей при изменении магнитного
потока $ \Phi $ в катушке, намотанной на образец. Пусть катушка c $ N $ витками плотно охватывает образец сечением $ S $, и индукция $ B $ в образце
однородна. Тогда

\begin{equation}
|B|=\frac{1}{SN}\int\mathcal{E} dt.
\label{eq:|B|}
\end{equation}

\begin{figure}[h]
\centering
    \includegraphics[width=0.3\linewidth]{3.4.5 Петля гистерезиса/pictures/gist2.jpg}
    \caption{Интегрирующая ячейка}
	\label{fig:int}
\end{figure}

Для интегрирования в работе используется интегрирующая $ RC $-цепочка (рис. $ \ref{fig:int} $).
«Входное» напряжение от источника $U_{\text{вх}}(t)$ подаётся на последовательно соединённые резистор $R_\text{и}$ и конденсатор $C_\text{и}$. <<Выходное>>
напряжение $U_{\text{вых}}(t)$ снимается с конденсатора. Предположим, что 1) сопротивление источника мало по сравнению с $R_\text{и}$, 2) выходное сопротивление (сопротивление на входе осциллографа), напротив, велико: $R_{\text{вых}}$ $ \gg $ $R_\text{и}$ и, наконец, 3) сопротивление $R_\text{и}$ достаточно велико, так что почти всё падение напряжения приходится на него, а $U_{\text{вых}}$ $\ll$ $U_{\text{вх}}$. В таком случае ток цепи равен I = ($U_{\text{вх}}$ - $U_{\text{вых}}$)/$R_\text{и}$ $\approx$ $U_{\text{вх}}$/$R_\text{и}$, и входное и выходное сопротивление связаны соотношением
\begin{equation}
U_{\text{вых}} = \frac{q}{C_\text{и}} = \frac{1}{C_\text{и}}\int\limits_0^t Idt \approx \frac{1}{\tau_\text{и}} \int\limits_0^t U_{\text{вх}}dt,
\label{eq:U_ext}
\end{equation}
где $\tau_\text{и}=R_\text{и}C_\text{и}$ - постоянная времени $ RC $ - цепочки. Для индукции поля из ($\ref{eq:|B|}$) получаем 
\begin{equation}
|B|=\frac{1}{SN}\int U_{\text{вх}} dt=\frac{\tau_\text{и}}{SN}U_{\text{вых}}.
\label{eq:|B|new}
\end{equation}

\textbf{Замечание.} Уточним критерий применимости соотношения ($\ref{eq:U_ext}$). Пусть на вход интегрирующей ячейки подан синусоидальный сигнал с частотой $\omega_0$. Тогда, пользуясь методом комплексных амплитуд, нетрудно найти отношение амплитуд входного и выходного напряжений:
\begin{equation}
\frac{U_{\text{вых}}}{U_{\text{вх}}}=\frac{1/\omega_0C}{\sqrt{R^2+1/(\omega_0C)^2}}.
\end{equation}
Тогда неравенство $U_{\text{вых}} \ll U_{\text{вх}}$ реализуется, если 
\begin{equation}
\tau \equiv RC\gg \frac{1}{\omega_0}
\end{equation}
(импеданс конденсатора мал по сравнению сопротивлением резистора).
В таком случае для синусоидального сигнала имеем
\begin{equation}
\frac{U_{\text{вых}}}{U_{\text{вх}}}\approx\frac{1}{\omega_0\tau}.
\end{equation}
В общем случае, если $\omega_0$ — частота самой низкой гармоники в спектре
произвольного входного сигнала, то при $\omega_0\tau \gg 1$ неравенство $U_{\text{вых}} \ll U_{\text{вх}}$ выполняется на любой частоте $\omega > \omega_0$.


\section*{Экспериментальная установка}

Схема экспериментальной установки показана на рис. 3.
	
Действующее значение переменного тока в обмотке $N_0$ измеряется амперметром А (мультиметром GDM). Последовательно с амперметром включено сопротивление $R_{0}$, напряжение с которого подается на вход X электронного осциллографа (ЭО). Это напряжение пропорционально току в обмотке $N_{0}$, а следовательно и напряженности H магнитного поля в образце.
	
Для измерения магнитной индукции B с измерительной обмотки $N_{И}$ на вход интегрирующей RC -цепочки подается напряжение $U_{\text{И}}$ ($U_{\text{ВХ}}$), пропорциональное производной $\dot{B}$, а с выхода снимается напряжение $U_{\text{C}}$($U_{\text{ВЫХ}}$), пропорциональное
величине B , и подается на вход Y осциллограа.
Замкнутая кривая, возникающая на экране, воспроизводит в некотором масштабе (различном для осей X и Y ) петлю гистерезиса. Чтобы придать этой кривой количественный смысл, необходимо установить масштабы изображения, т.е. провести калибровку каналов X и Y ЭО. Для этого, впервых, надо узнать, каким напряжениям (или токам) соответствуют амплитуды сигналов, видимых на экране, и во-вторых,  каким значениям B и H соответствуют эти напряжения
(или токи).

\begin{figure}[h!]
	\centering
	\includegraphics[width=0.8\linewidth]{3.4.5 Петля гистерезиса/pictures/gist.jpg}
	\caption{ Схема установки для исследования намагничивания образцов}
	\label{fig:scheme}
\end{figure}

\section*{Задание}
\begin{center}
    I. Петля гистерезиса на экране ЭО
\end{center}

1) схема собрана как на рисунке 3 и осциллограф настроен под установку по инструкции

2) 
\begin{figure}[h]
\centering
\includegraphics[width=0.3\linewidth]{3.4.5 Петля гистерезиса/pictures/петля пермаллой.png}
\label{fig:mpr}
\end{figure}

Сопротивление исчезновения усов: $R^{min}_{\text{пот}} = \frac{1}{7} R_{max}$

Сопротивление с максимальной площадью: $R^{max}_{\text{пот}} = \frac{3}{7} R_{max}$

3) - 6) Подбираем коефициенты для максимальной занимаемой площади и определяем параметры первого образца, пермаллой:

\textit{Параметры:} 
\begin{itemize}
    \item $N_0 = 15$
    \item $N_{\text{и}} = 300$
    \item $S = 0,66$ см$^2$
    \item $2\pi R = 14,1$ см
    
\end{itemize}

\begin{equation*}
    \begin{cases}
        K_X = 20 \;\text{мВ}
        \\
        K_Y = 50 \;\text{мВ}
    \end{cases}
    ,
    \begin{cases}
        2\cdot X_S = 176 \;\text{мВ}
        \\
        2\cdot X_Y = 180 \;\text{мВ}
    \end{cases}
    ,
    \begin{cases}
        [2\cdot X_S] = 120 \;\text{мВ}
        \\
        [2\cdot X_Y] = 160 \;\text{мВ}
    \end{cases}
\end{equation*}

\begin{table}[h!]
  		\caption{Начальная кривая намагничивания для первого образца, пермаллоя}
  		\begin{center}
  			\begin{tabular}{|c|c|c|c|c|c|c|c|} 
  				\hline 
  				№ &  1 &  2 & 3 & 4 & 5 &  6  & 7 \\ 	\hline
  				
  			 $ x $, дел & 4.9 & 3.9 & 3.6 & 2.8 & 2.0 & 1.5 & 0.0 \\
  				 $ y $, дел & 1.8 & 1.4 & 1.1 & 0.7  & 0.35 & 0.2 & 0.0 \\
  				\hline
  				
  			\end{tabular}
  		\end{center}
  	\end{table}

\begin{figure}[h]
\centering
\includegraphics[width=0.7\linewidth]{3.4.5 Петля гистерезиса/pictures/Пермаллой.png}
\end{figure}

Повторяем тоже самое для двух других образцов: 8) \\
\textit{Феррит}

\begin{figure}[h]
\centering
\includegraphics[width=0.3\linewidth]{3.4.5 Петля гистерезиса/pictures/петля Феррит.jpg}
\label{fig:mpr}
\end{figure}

\textit{Параметры:} 
\begin{itemize}
    \item $N_0 = 42$
    \item $N_{\text{и}} = 400$
    \item $S = 3,0$ см$^2$
    \item $2\pi R = 25,0$ см
    
\end{itemize}

\begin{equation*}
    \begin{cases}
        K_X = 50 \;\text{мВ}
        \\
        K_Y = 50 \;\text{мВ}
    \end{cases}
    ,
    \begin{cases}
        2\cdot X_S = 220 \;\text{мВ}
        \\
        2\cdot X_Y = 110 \;\text{мВ}
    \end{cases}
    ,
    \begin{cases}
        [2\cdot X_S] = 60 \;\text{мВ}
        \\
        [2\cdot X_Y] = 50 \;\text{мВ}
    \end{cases}
\end{equation*}

\begin{table}[h!]
  		\caption{Начальная кривая намагничивания Феррита}
  		\begin{center}
  			\begin{tabular}{|c|c|c|c|c|c|c|c|} 
  				\hline 
  				№ &  1 &  2 & 3 & 4 & 5 &  6  & 7 \\ 	\hline
  				
  			 $ x $, дел & 2.2 & 1.8 & 1.4 & 1.2 & 1.0 & 0.8 & 0.0 \\
  				 $ y $, дел & 1.1 & 1.0 & 0.8 & 0.7 & 0.6 & 0.4 & 0.0 \\
  				\hline
  				
  			\end{tabular}
  		\end{center}
  	\end{table}
\begin{figure}[h]
\centering
\includegraphics[width=0.68\linewidth]{3.4.5 Петля гистерезиса/pictures/Феррин.pdf}
\end{figure} \newpage

\textit{Кремнистое железо}

\begin{figure}[h]
\centering
\includegraphics[width=0.2\linewidth]{3.4.5 Петля гистерезиса/pictures/петля Кремнистое железо.jpg}
\label{fig:mpr}
\end{figure}

\textit{Параметры:} 
\begin{itemize}
    \item $N_0 = 20$
    \item $N_{\text{и}} = 200$
    \item $S = 2,0$ см$^2$
    \item $2\pi R = 11,0$ см
    
\end{itemize}

\begin{equation*}
    \begin{cases}
        K_X = 100 \;\text{мВ}
        \\
        K_Y = 20 \;\text{мВ}
    \end{cases}
    ,
    \begin{cases}
        2\cdot X_S = 680 \;\text{мВ}
        \\
        2\cdot X_Y = 140 \;\text{мВ}
    \end{cases}
    ,
    \begin{cases}
        [2\cdot X_S] = 160 \;\text{мВ}
        \\
        [2\cdot X_Y] = 136 \;\text{мВ}
    \end{cases}
\end{equation*}

\begin{table}[h!]
  		\caption{Начальная кривая намагничивания для кремнистого железа}
  		\begin{center}
  			\begin{tabular}{|c|c|c|c|c|c|c|c|} 
  				\hline 
  				№ &  1 &  2 & 3 & 4 & 5 &  6  & 7 \\ 	\hline
  				
  			 $ x $, дел & 3.4 & 2.4 & 1.7 & 1.2 & 1.0 & 0.6 & 0.0 \\
  				 $ y $, дел & 3,5 & 3.1 & 2.7 & 2.1  & 1.6 & 0.7 & 0.0 \\
  				\hline
  				
  			\end{tabular}
  		\end{center}
  	\end{table}

\begin{figure}[h]
\centering
\includegraphics[width=0.7\linewidth]{3.4.5 Петля гистерезиса/pictures/Кремнеевое железо.pdf}
\end{figure}

\begin{center}
    II. Проверка калибровки осциллографа
\end{center}
9) проводим калибровку по горизонтальной оси закорачивая обмотку на $N_0$:
\begin{itemize}
    \item 
    \begin{equation*}
    \begin{cases}
        K_X = 0.1 \;\text{В}
        \\
        I_1 = 1.4 \;\text{А}
        \\
        2 X_1 = 9.8 
    \end{cases}
    \Rightarrow 
    K_{X_1}^{\text{теор}} = 2R_0 \sqrt{2} I_{\text{эф}}/(2X_1) = 0.08 \;\text{В/дел}
    \end{equation*}

    \begin{equation*}
    \begin{cases}
        K_X = 0.1 \;\text{В}
        \\
        I_2 = 1.08 \;\text{А}
        \\
        2 X_2 = 7.8 
    \end{cases}
    \Rightarrow 
    K_{X_2}^{\text{теор}} = 2R_0 \sqrt{2} I_{\text{эф}}/(2X_1) = 0.78 \;\text{В/дел}
    \end{equation*}
    Получаем отличие от получаемого на экране $\sim20\%$
    \item 
    \begin{equation*}
    \begin{cases}
        K_X = 50 \;\text{мВ}
        \\
        I_1 = 0.68 \;\text{А}
        \\
        2 X_1 = 9.8 
    \end{cases}
    \Rightarrow 
    K_{X_1}^{\text{теор}} = 2R_0 \sqrt{2} I_{\text{эф}}/(2X_1) = 39.3 \;\text{мВ/дел}
    \end{equation*}

    \begin{equation*}
    \begin{cases}
        K_X = 50 \;\text{мВ}
        \\
        I_2 = 1.03 \;\text{А}
        \\
        2 X_2 = 7.4 
    \end{cases}
    \Rightarrow 
    K_{X_2}^{\text{теор}} = 2R_0 \sqrt{2} I_{\text{эф}}/(2X_1) = 39.5 \;\text{мВ/дел}
    \end{equation*}
    Получаем отличие от получаемого на экране $\sim20\%$
    \item 
    \begin{equation*}
    \begin{cases}
        K_X = 20 \;\text{мВ}
        \\
        I_1 = 0.256 \;\text{А}
        \\
        2 X_1 = 9.2 
    \end{cases}
    \Rightarrow 
    K_{X_1}^{\text{теор}} = 2R_0 \sqrt{2} I_{\text{эф}}/(2X_1) = 16 \;\text{мВ/дел}
    \end{equation*}

    \begin{equation*}
    \begin{cases}
        K_X = 20 \;\text{мВ}
        \\
        I_2 = 0.160 \;\text{А}
        \\
        2 X_2 = 6.4 
    \end{cases}
    \Rightarrow 
    K_{X_2}^{\text{теор}} = 2R_0 \sqrt{2} I_{\text{эф}}/(2X_1) = 14 \;\text{мВ/дел}
    \end{equation*}
    Получаем отличие от получаемого на экране $\sim20\%$, но в последнем измерении была допущена измерительная ошибка (стало ясно потом), измерение не учитывается
\end{itemize}

Погрешности получал по формуле: $\Delta K^{\text{теор}} = \left( \frac{\Delta R_0}{R_0} + \frac{\Delta I_{\text{эф}}}{I_{\text{эф}}} + \frac{\Delta X_i}{X_i}\right) \cdot K_{\text{теор}}$, где погрешности $\Delta R_0 \rightarrow 0$, $\Delta I_{\text{эф}} = 0.003$ А - определется приборной погрешностью, $\Delta X_i = 0.1$  - определется приборной погрешностью. \\ 
На этом мое выполнение было закончено 10) - 13) пропускаем. \\
14) будем пользоваться следующими формулами для определения масштаба: 
\begin{equation*}
    I = \frac{I N_0}{2 \pi R} = \frac{K_S N_0}{R_0 2\pi R}
\end{equation*}
\begin{equation*}
    B = \frac{R_{\text{и}} U_{\text{и}}}{S N_{\text{и}}} K_Y
\end{equation*}

\begin{center}
    \begin{tabular}{|l|r|r|r|}
        \hline
        образец & \multicolumn{1}{l|}{Пармаллоид} & \multicolumn{1}{l|}{Феррит} & \multicolumn{1}{l|}{Кремнистое железо} \\ \hline
        $H$, $\left( \frac{\text{А}}{\text{м}} \right)$ / дел       & 10.6                            & 42.0                        & 90.9                                   \\ \hline
        $B$, Тл / дел       & 1.0                             & 0.2                         & 0.2                                    \\ \hline
    \end{tabular}
\end{center}
15) Делаем вычисления на основе получившихся картин на экране осциллографа: 
\begin{center}
    \begin{tabular}{|l|r|r|r|}
        \hline
        образец & \multicolumn{1}{l|}{Пармаллоид} & \multicolumn{1}{l|}{Феррит} & \multicolumn{1}{l|}{Кремнистое железо} \\ \hline
        $H_{max}$, $\left( \frac{\text{А}}{\text{м}} \right)$      & 17.0                            & 96.6                        & 318.2                                   \\ \hline
        $B_S$, Тл       & 1.8                             & 0.2                         & 0.7                                    \\ \hline
    \end{tabular}
\end{center}
16) Считаем коэрцитивное поле $H_c$ и остаточную индукцию $B_r$
\begin{center}
    \begin{tabular}{|l|r|r|r|}
        \hline
        образец & \multicolumn{1}{l|}{Пармаллоид} & \multicolumn{1}{l|}{Феррит} & \multicolumn{1}{l|}{Кремнистое железо} \\ \hline
        $H_{c}$, $\left( \frac{\text{А}}{\text{м}} \right)$      & 14.8                            & 29.6                        & 81.8                                   \\ \hline
        $B_r$, Тл       & 1.6                             & 0.1                         & 0.3                                    \\ \hline
    \end{tabular}
\end{center}
17) Считаем $\mu_{\text{диф}} = \frac{dB}{dH}$:
\begin{center}
    \begin{tabular}{|l|r|r|r|}
        \hline
        образец & \multicolumn{1}{l|}{Пармаллоид} & \multicolumn{1}{l|}{Феррит} & \multicolumn{1}{l|}{Кремнистое железо} \\ \hline
        $\mu_{\text{диф}}$, $\frac{\text{дел}}{\text{дел}}$ & 0.1 & 0.4 & 1 \\ \hline
        $\mu_{\text{диф}}$,  $\left( \frac{\text{А}}{\text{м}\cdot\text{Тл}} \right)\cdot 10^{-3}$     & 9.4                            & 1.9                        & 2.2                                   \\ \hline
    \end{tabular}
\end{center}
Погрешности по нашим данным оценить тяжело, поэтому будет честно взять погрешность $\sim 50 \%$

\section*{Вывод}
Итого мы с помощью динамическог метода основанного на свойствах интегрирующей цепочки смогли наблюдать кривые гистерезиса для выданных нам образцов, а также измерить коэрцетивную напряженность, остаточную индукцию и индукцию с напряженностью насыщения. Отличия от табличных значений в среднем на 15\%.
Также мы прокалибровали ОЭ по оси X и получили что значения отличаются в районе 20\%.

\end{document}
